\section{Neutron Star Structure}
\label{sec:ns_structure}

This section documents Module~1 of the visualization pipeline: the computation
of static, spherically symmetric neutron star models via the
Tolman--Oppenheimer--Volkoff (TOV) equations, closed by the Brussels--Montreal
BSk unified equations of state.  Every equation implemented in the code is
stated below, cross-referenced against the source literature, and any
approximations or discrepancies are flagged.

%==========================================================================
\subsection{Physical constants}
\label{sec:constants}

All physical constants are taken in CGS units from the CODATA 2018
adjustment~\cite{Tiesinga2021}.  The values used are:
%
\begin{equation}
  c = 2.997\,924\,58\times10^{10}\;\text{cm\,s}^{-1},\qquad
  G = 6.674\,30\times10^{-8}\;\text{cm}^{3}\,\text{g}^{-1}\,\text{s}^{-2}.
  \label{eq:cG}
\end{equation}
%
The solar mass is set to
$M_{\odot}=1.988\,92\times10^{33}\;\text{g}$, attributed to IAU 2015 nominal
values.

\paragraph{Remark.}
The IAU 2015 nominal solar mass parameter is
$\mathcal{GM}_{\odot}^{N}=1.327\,124\,4\times10^{26}\;\text{cm}^{3}\,
\text{s}^{-2}$~\cite{IAU2015},
which combined with the CODATA 2018 value of $G$ yields
$M_{\odot}\approx1.988\,48\times10^{33}\;\text{g}$.
The code value $1.988\,92\times10^{33}\;\text{g}$ is a commonly used older
estimate and differs by ${\sim}0.02\%$.  This has negligible impact on
computed stellar masses at the precision of the EOS fit (${\sim}1\%$).

%==========================================================================
\subsection{BSk equation of state}
\label{sec:bsk_eos}

The equation of state (EOS) is an analytical representation of the unified
Brussels--Montreal BSk models, following
Potekhin et al.\ (2013)~\cite{Potekhin2013} (hereafter P13).
The code implements three parameter sets---BSk19, BSk20, and BSk21---with
BSk21 used as the default (stiffest EOS, maximum mass
$M_{\max}\approx2.27\,M_{\odot}$).

\subsubsection{Pressure--density relation}

Introducing the dimensionless logarithmic variables
\begin{equation}
  \xi \equiv \log_{10}\!\bigl(\rho/\text{g\,cm}^{-3}\bigr),\qquad
  \zeta \equiv \log_{10}\!\bigl(P/\text{dyn\,cm}^{-2}\bigr),
  \label{eq:xi_zeta_def}
\end{equation}
the analytical fit for $P(\rho)$ is given by P13, Eq.~(3):
%
\begin{align}
  \zeta &= \frac{a_1 + a_2\,\xi + a_3\,\xi^{3}}{1 + a_4\,\xi}\;
           \bigl\{\exp\!\bigl[a_5(\xi - a_6)\bigr]+1\bigr\}^{-1}
           \notag\\[4pt]
  &\quad + (a_7 + a_8\,\xi)\;
           \bigl\{\exp\!\bigl[a_9(a_6 - \xi)\bigr]+1\bigr\}^{-1}
           \notag\\[4pt]
  &\quad + (a_{10} + a_{11}\,\xi)\;
           \bigl\{\exp\!\bigl[a_{12}(a_{13} - \xi)\bigr]+1\bigr\}^{-1}
           \notag\\[4pt]
  &\quad + (a_{14} + a_{15}\,\xi)\;
           \bigl\{\exp\!\bigl[a_{16}(a_{17} - \xi)\bigr]+1\bigr\}^{-1}
           \notag\\[4pt]
  &\quad + \frac{a_{18}}{1+\bigl[a_{19}(\xi - a_{20})\bigr]^{2}}
         + \frac{a_{21}}{1+\bigl[a_{22}(\xi - a_{23})\bigr]^{2}}\,.
  \label{eq:bsk_pressure}
\end{align}
%
The pressure is then $P = 10^{\zeta}$.

\paragraph{Verification against paper.}
The implementation in \texttt{bsk\_eos.jl} (lines 121--127) reproduces
Eq.~\eqref{eq:bsk_pressure} term-by-term with no discrepancies.
Each of the six additive terms matches P13 Eq.~(3) identically:
four Fermi--Dirac-type terms and two Lorentzian terms.

\paragraph{Validity range.}
The fit is valid for
$10^{6}\lesssim\rho\lesssim10^{16}\;\text{g\,cm}^{-3}$
($6\lesssim\xi\lesssim16$).
The code enforces $\rho>0$ and $\rho<10^{17}\;\text{g\,cm}^{-3}$ via
assertions.  Per P13, the typical fit error on $P$ is ${\approx}\,1\%$ for
$\xi\gtrsim6$, with a maximum of $3.7\%$ at $\xi=9.51$ (the neutron-drip
transition region).

\subsubsection{Fit parameters}

The 23 parameters $a_1,\ldots,a_{23}$ for each BSk model are taken from
P13 Table~2.  All parameter values in the code have been checked against the
published table; the comparison for the preferred BSk21 model is shown in
Table~\ref{tab:bsk21_params}.

\begin{table}[htbp]
\centering
\caption{BSk21 parameters $a_i$: code values versus P13 Table~2.
All entries match exactly to the published precision.}
\label{tab:bsk21_params}
\begin{tabular}{r r r}
\hline
$i$ & Code & P13 Table~2 \\
\hline
 1 & $4.857$    & $4.857$ \\
 2 & $6.981$    & $6.981$ \\
 3 & $0.00706$  & $0.00706$ \\
 4 & $0.19351$  & $0.19351$ \\
 5 & $4.085$    & $4.085$ \\
 6 & $12.065$   & $12.065$ \\
 7 & $10.521$   & $10.521$ \\
 8 & $1.5905$   & $1.5905$ \\
 9 & $4.104$    & $4.104$ \\
10 & $-28.726$  & $-28.726$ \\
11 & $2.0845$   & $2.0845$ \\
12 & $4.89$     & $4.89$ \\
13 & $14.302$   & $14.302$ \\
14 & $22.881$   & $22.881$ \\
15 & $-1.7690$  & $-1.7690$ \\
16 & $0.989$    & $0.989$ \\
17 & $15.313$   & $15.313$ \\
18 & $0.091$    & $0.091$ \\
19 & $4.68$     & $4.68$ \\
20 & $11.65$    & $11.65$ \\
21 & $-0.086$   & $-0.086$ \\
22 & $10.0$     & $10.0$ \\
23 & $14.15$    & $14.15$ \\
\hline
\end{tabular}
\end{table}

The same check was performed for BSk19 and BSk20.
All $3\times23=69$ parameter values match the published table to every
quoted digit.  No transcription errors were found.

\subsubsection{Energy density}

For the TOV system, the total energy density $\varepsilon$ is required.
Following P13 Sect.~3.1 (discussion below their Eq.~1), the variable $\rho$
in the BSk parametrization is the \emph{gravitational mass density}, which
already incorporates internal (binding) energy contributions.  The
code therefore uses
\begin{equation}
  \varepsilon = \rho\,c^{2}\,,
  \label{eq:energy_density}
\end{equation}
which is the correct identification for this parametrization.

\paragraph{Approximation.}
This relation is exact within the P13 fitting framework.
In a more fundamental treatment, one would distinguish rest-mass energy
density $\rho_0 c^{2}$ from thermal and interaction contributions.
The BSk fit absorbs these corrections into the definition of $\rho$, so
Eq.~\eqref{eq:energy_density} is self-consistent.

\subsubsection{Pseudo-enthalpy}
\label{sec:pseudo_enthalpy}

The pseudo-enthalpy, useful for rotating-star calculations, is defined by
P13 Eq.~(2):
\begin{equation}
  H = \int_{0}^{P}\frac{\mathrm{d}P'}{\rho(P')\,c^{2}+P'}\,.
  \label{eq:pseudo_enthalpy}
\end{equation}
The code evaluates this integral via the trapezoidal rule in $\log_{10}P$
space with $N=1000$ points (default), using the change of variables
$\mathrm{d}P' = P'\ln10\;\mathrm{d}(\log_{10}P')$, so that the integrand
becomes
\begin{equation}
  H = \int_{\log_{10}P_{\min}}^{\log_{10}P}
      \frac{P'\,\ln 10}{\varepsilon(P') + P'}\;\mathrm{d}(\log_{10}P')\,,
  \label{eq:pseudo_enthalpy_logP}
\end{equation}
where $P_{\min}=P(\rho=10^{6}\;\text{g\,cm}^{-3})$ is the lower limit of
the EOS fit validity.

\paragraph{Verification against paper.}
Equation~\eqref{eq:pseudo_enthalpy} matches P13 Eq.~(2) identically.
The logarithmic integration variable is a numerical convenience not
present in the paper; it improves quadrature accuracy over the many
decades of pressure spanning the stellar interior.

\paragraph{Approximation.}
The function $\rho(P)$ required inside the integrand is obtained by
bisection inversion of $P(\rho)$ (see Eq.~\ref{eq:bsk_pressure}), with
relative tolerance $10^{-10}$ in $\log_{10}\rho$.  The trapezoidal rule with
$N=1000$ uniform steps in $\log_{10}P$ yields accuracy well within the
underlying ${\sim}1\%$ EOS fit error.

\subsubsection{Density inversion}

Since the BSk fit provides $P(\rho)$ analytically but $\rho(P)$ is needed
for the TOV integration, the inverse function is computed by bisection in
$\log_{10}\rho$ space:
\begin{equation}
  \text{find } \rho^{*} \text{ such that }
  P(\rho^{*}) = P_{\text{target}},
  \quad \rho^{*}\in[10^{4},\,10^{16}]\;\text{g\,cm}^{-3}.
  \label{eq:bisection}
\end{equation}
The bracket is automatically extended by factors of~10 if the target
pressure falls outside the initial range.  Convergence to relative
tolerance $10^{-10}$ is guaranteed within 200 iterations
($2^{200}\gg10^{12}$, the bracket ratio in decades).

\paragraph{Approximation.}
Bisection is unconditionally stable but slower than Newton--Raphson.
Since $\mathrm{d}P/\mathrm{d}\rho$ is not available analytically from
the fit (it would require differentiating
Eq.~\ref{eq:bsk_pressure}), bisection is the simplest robust choice.
The ${\sim}40$ iterations typically needed add negligible cost compared to
the RK4 TOV loop.

%==========================================================================
\subsection{TOV equations}
\label{sec:tov}

The structure of a static, spherically symmetric star in general relativity is
governed by the TOV equations
\cite{Tolman1939,OppenheimerVolkoff1939,ShapiroTeukolsky1983}.
Following the notation of Shapiro \& Teukolsky (1983), Ch.~5--6:

\begin{align}
  \frac{\mathrm{d}P}{\mathrm{d}r}
  &= -\,\frac{G\,(\varepsilon + P)\!
       \left(m + \dfrac{4\pi\,r^{3}\,P}{c^{2}}\right)}
       {r\!\left(r\,c^{2} - 2\,G\,m\right)}\,,
  \label{eq:tov_dPdr}\\[8pt]
  \frac{\mathrm{d}m}{\mathrm{d}r}
  &= \frac{4\pi\,r^{2}\,\varepsilon}{c^{2}}\,,
  \label{eq:tov_dmdr}
\end{align}
%
where $r$ is the circumferential (areal) radius, $m(r)$ is the enclosed
gravitational mass, $P(r)$ is the pressure, and
$\varepsilon(r)=\rho(r)\,c^{2}$ is the total energy density.  The system is
closed by the EOS $P=P(\rho)$ from Sect.~\ref{sec:bsk_eos}.

\paragraph{Verification against paper.}
The code implementation in \texttt{tov.jl} (lines 167--175) reproduces
Eqs.~\eqref{eq:tov_dPdr}--\eqref{eq:tov_dmdr} exactly.  The function
\texttt{tov\_rhs} returns $(dm/dr,\;dP/dr)$ with the signs and grouping
matching the standard form.  An assertion checks that the denominator
$r(rc^{2}-2Gm)>0$, i.e.\ that the integration has not crossed the
Schwarzschild horizon.

\subsubsection{Central boundary conditions}
\label{sec:central_bc}

At $r=0$ the TOV equations are singular ($0/0$ form).  The code avoids
this by starting the integration at a small radius $r_{0}$ using a
Taylor expansion about the centre.  To derive these, one substitutes
$m(r)=(4\pi/3)(\varepsilon_{c}/c^{2})\,r^{3}+\mathcal{O}(r^{5})$ into
Eq.~\eqref{eq:tov_dPdr} and expands to leading non-trivial order:
%
\begin{align}
  m(r_0) &= \frac{4\pi}{3}\,\frac{\varepsilon_{c}}{c^{2}}\,r_0^{3}\,,
  \label{eq:m_central}\\[6pt]
  P(r_0) &= P_{c} - \frac{2\pi\,G}{3\,c^{4}}\,
             (\varepsilon_{c}+P_{c})(\varepsilon_{c}+3P_{c})\,r_0^{2}\,.
  \label{eq:P_central}
\end{align}
%
These are the standard results (see, e.g., Shapiro \& Teukolsky 1983,
Eq.~(5.7.6); Lindblom 1992~\cite{Lindblom1992}).

\paragraph{Verification against code.}
The code (\texttt{tov.jl}, lines 73--75) implements:
\begin{verbatim}
  m0 = 4pi/3 * eps_c / c^2 * r0^3
  dP_coeff = -2pi * G * (eps_c + P_c) * (eps_c + 3P_c) / (3 * c^4)
  P0 = P_c + dP_coeff * r0^2
\end{verbatim}
This matches Eqs.~\eqref{eq:m_central}--\eqref{eq:P_central} exactly.
The default starting radius is $r_0 = 100\;\text{cm} = 1\;\text{m}$, which
is negligible compared to stellar radii of order $10^{6}\;\text{cm}$ and
ensures the $\mathcal{O}(r^{4})$ truncation error is of order
$(r_{0}/R)^{4}\sim10^{-16}$.

\subsubsection{Numerical integration}

The TOV system is integrated outward using a classical fourth-order
Runge--Kutta (RK4) scheme.  The step size $\Delta r$ is adaptive:
\begin{equation}
  \Delta r = \min\!\Bigl(
    \Delta r_{\max},\;
    \Delta r_{\text{init}}\cdot\max\!\bigl(1,\,(P/P_{\text{surf}})^{0.3}\bigr),\;
    0.01\,r
  \Bigr),
  \label{eq:adaptive_step}
\end{equation}
with defaults $\Delta r_{\text{init}}=100\;\text{cm}$,
$\Delta r_{\max}=1000\;\text{cm}$.
The three constraints ensure that steps are (i) bounded above, (ii) small
near the surface where $P$ changes rapidly, and (iii) never exceed $1\%$ of
the current radius.

\paragraph{Approximation.}
The step controller in Eq.~\eqref{eq:adaptive_step} is heuristic rather
than based on local truncation error estimation.  It does not use
embedded Runge--Kutta pairs (e.g.\ Dormand--Prince) for rigorous error
control.  However, the small maximum step size
($\Delta r_{\max}=10\;\text{m}$, compared to $R\sim10\;\text{km}$) and the
additional $1\%$-of-$r$ constraint ensure that the integration is well
resolved.

\paragraph{Surface termination.}
Integration terminates when $P$ drops below the surface pressure
$P_{\text{surf}} \equiv P(\rho=10^{6}\;\text{g\,cm}^{-3})$, which
corresponds to the lower validity limit of the BSk fit.  At the final step,
a linear interpolation in $r$ locates the exact surface crossing:
\begin{equation}
  r_{\text{surf}} = r_{n-1}
    + \frac{P_{n-1}-P_{\text{surf}}}{P_{n-1}-P_{n}}\;\Delta r\,,
  \label{eq:surface_interp}
\end{equation}
and the mass is interpolated likewise.

\subsubsection{Derived quantities}

From the integration output---the total gravitational mass $M\equiv m(R)$
and circumferential radius $R\equiv r_{\text{surf}}$---the code computes
the following:

\paragraph{Compactness.}
\begin{equation}
  u \equiv \frac{2\,G\,M}{R\,c^{2}}\,.
  \label{eq:compactness}
\end{equation}
This is the standard relativistic compactness parameter (the ratio of
Schwarzschild radius to stellar radius).

\paragraph{Surface gravity.}
\begin{equation}
  g_{\text{surf}} = \frac{G\,M}{R^{2}\,\sqrt{1-u}}
                   = \frac{G\,M}{R^{2}}\,
                     \left(1 - \frac{2\,G\,M}{R\,c^{2}}\right)^{\!-1/2}.
  \label{eq:surface_gravity}
\end{equation}
%
This is the \emph{local} surface gravity, \emph{not} the
redshifted gravity $g_{\infty}=g_{\text{surf}}\sqrt{1-u}=GM/R^{2}$
seen by a distant observer.  The code comment labels it
``redshifted'', but the formula actually gives the local value---the
redshifted surface gravity would instead be $g_{\infty}=GM/R^{2}$.

\paragraph{Discrepancy.}
\emph{The code comment on line 184 of \texttt{tov.jl} describes
\texttt{g\_surf} as ``surface gravity (redshifted)'', but the implemented
formula $g = GM/\bigl(R^{2}\sqrt{1-2GM/Rc^{2}}\bigr)$ is the local
(proper) surface gravity, not the redshifted one.  The formula itself is
correct for the local surface gravity; only the comment is misleading.}

\subsubsection{Mass--radius curve}

The function \texttt{mass\_radius\_curve} varies the central density
logarithmically over
$\log_{10}(\rho_{c}/\text{g\,cm}^{-3})\in[14.0,\,15.8]$ and calls
\texttt{solve\_tov} at each value to trace the $M(R)$ relation.
This range spans from $\rho_{c}\approx10^{14}\;\text{g\,cm}^{-3}$
(slightly below nuclear saturation) to
$\rho_{c}\approx6.3\times10^{15}\;\text{g\,cm}^{-3}$
(well above the maximum-mass configuration).

%==========================================================================
\subsection{Summary of approximations}
\label{sec:approx_summary}

The key approximations and simplifications in Module~1 are:
\begin{enumerate}
  \item \textbf{EOS fit accuracy:} The BSk analytical fit
    (Eq.~\ref{eq:bsk_pressure}) reproduces the tabulated EOS to
    ${\sim}1\%$ for $\xi\gtrsim6$, with a worst case of $3.7\%$ at
    $\xi=9.51$ (P13, Sect.~3.2).
  \item \textbf{Zero temperature:} The BSk parametrization assumes cold,
    catalyzed matter ($T=0$).  Thermal corrections are neglected.
  \item \textbf{Spherical symmetry:} Rotation is neglected; the star is
    treated as static and spherically symmetric.
  \item \textbf{Heuristic step control:} The RK4 integrator uses an
    empirical adaptive step (Eq.~\ref{eq:adaptive_step}) rather than
    embedded error estimation.
  \item \textbf{Solar mass value:} $M_{\odot}$ differs from the precise
    IAU~2015 nominal value by ${\sim}0.02\%$
    (Sect.~\ref{sec:constants}).
  \item \textbf{Density inversion by bisection:} $\rho(P)$ is obtained
    numerically rather than from a separate analytical inverse fit
    (P13 provide such fits in their Eqs.~6--7, which are not used here).
\end{enumerate}

\subsection{Summary of discrepancies}
\label{sec:discrepancies}

\begin{enumerate}
  \item \textbf{Surface gravity comment (cosmetic):} The code comment on
    \texttt{tov.jl} line~184 labels \texttt{g\_surf} as
    ``redshifted'', but the formula computes the local proper surface
    gravity.  The formula is correct; only the comment is wrong.
    See Eq.~\eqref{eq:surface_gravity}.
  \item \textbf{Solar mass value (minor):} $M_{\odot}=1.98892\times
    10^{33}\;\text{g}$ vs.\ the IAU~2015/CODATA~2018 derived value of
    $1.98848\times10^{33}\;\text{g}$ (a $0.02\%$ difference).
  \item \textbf{No discrepancies in physics equations:} All implemented
    equations---the BSk pressure fit (Eq.~\ref{eq:bsk_pressure}),
    TOV system (Eqs.~\ref{eq:tov_dPdr}--\ref{eq:tov_dmdr}), central
    expansion (Eqs.~\ref{eq:m_central}--\ref{eq:P_central}), and
    pseudo-enthalpy (Eq.~\ref{eq:pseudo_enthalpy})---match the source
    references exactly.  All 69 BSk fit parameters across three models
    match P13 Table~2 to every quoted digit.
\end{enumerate}
